%% ============================================================================
%%  back/publications.tex
%%  فهرست مقاله‌های مستخرج از پایان‌نامه / رساله
%%  طراحی شده برای قالب رسمی دانشگاه خوارزمی
%% ============================================================================
%%
%%  توضیح:
%%  در این بخش، مشخصات کتاب‌شناختی مقالات علمی-پژوهشی، مقالات کنفرانسی بین‌المللی
%%  یا اختراعات ثبت‌شده‌ای که حاصل یافته‌های مستقیم این پایان‌نامه/رساله هستند قید می‌گردد.
%% ============================================================================

\listofpublications{%
\section*{الف) مقالات چاپ‌شده در نشریات علمی-پژوهشی و بین‌المللی (مجلات)}
\begin{enumerate}[label={[\arabic*]}]
  \item آیریک پژوهشگر، دکتر احمد رضایی، و دکتر سارا محمدی،
  «طراحی و پیاده‌سازی قالب جامع و ماژولار حروف‌چینی دانشگاهی بر پایه موتور زی‌پرشین»،
  \textit{فصلنامه علمی-پژوهشی پردازش علائم و داده‌ها}،
  دوره ۲۱، شماره ۲، صفحات ۴۵--۶۲، تابستان ۱۴۰۳.
\end{enumerate}
\begin{latin}
\begin{enumerate}[label={[\arabic*]}, start=2]
  \item A. Researcher, A. Rezaei, and S. Mohammadi,
  ``Robust Bidirectional Document Formatting Engine for Scientific Theses Using XePersian,''
  \textit{IEEE Transactions on Professional Communication},
  vol.~67, no.~3, pp.~285--301, September 2024.
\end{enumerate}
\end{latin}

\vspace{0.8cm}
\section*{ب) مقالات ارائه‌شده در همایش‌ها و کنفرانس‌های علمی}
\begin{enumerate}[label={[\arabic*]}]
  \item آیریک پژوهشگر و دکتر احمد رضایی،
  «تحلیل مقایسه‌ای سیستم‌های خودکار حروف‌چینی دانشگاهی در زبان‌های راست‌به‌چپ»،
  \textit{سی‌ودومین کنفرانس ملی مهندسی برق و کامپیوتر ایران}،
  دانشگاه تهران، اردیبهشت ۱۴۰۳.
\end{enumerate}
\begin{latin}
\begin{enumerate}[label={[\arabic*]}, start=2]
  \item A. Researcher and A. Rezaei,
  ``Automating Academic Thesis Styles in Persian: Challenges and Solutions,''
  \textit{Proceedings of the 9th International Conference on Computer Science and Information Technology (CSIT)},
  Tehran, Iran, pp.~112--118, May 2024.
\end{enumerate}
\end{latin}
}
